\chapter{Segitiga dan Segi Empat}\label{ch:g6:shapes}

Segitiga dan \omterm{def:g4:geometry:polygon}{segi empat} adalah bangun dasar dalam geometri. Tujuan bab
ini adalah mengenal keluarga istimewanya --- \omterm{def:g6:shapes:triangles}{segitiga sama kaki}, sama
sisi, siku-siku; persegi panjang, \omterm{def:g6:shapes:quadrilaterals}{belah ketupat}, persegi --- lewat
\emph{sifat yang mendefinisikannya}, lalu melukisnya dengan tepat
memakai penggaris dan jangka.

\section{Segitiga}

\begin{definition}[Segitiga istimewa]\label{def:g6:shapes:triangles}
Segitiga $ABC$ punya tiga \omterm{def:g5:solids:def}{titik sudut} dan tiga sisi. Ia disebut:
\begin{itemize}
  \item \emph{sama kaki}\index{segitiga sama kaki} di $A$ jika
        $AB = AC$ (dua sisi sama panjang);
  \item \emph{sama sisi}\index{segitiga sama sisi} jika
        $AB = BC = CA$ (tiga sisi sama panjang);
  \item \emph{siku-siku}\index{segitiga siku-siku} di $A$ jika sudut
        $\widehat{BAC}$ adalah \omterm{def:g3:shapes:rightangle}{sudut siku-siku}. Sisi $[BC]$ yang
        berhadapan dengan \omterm{def:g3:shapes:rightangle}{sudut siku-siku} itu disebut \emph{sisi miring}.
\end{itemize}
\end{definition}

\begin{omfigure}
\begin{tikzpicture}[scale=0.85]
  % sama kaki
  \draw[thick] (0,0) -- (2.4,0) -- (1.2,2.1) -- cycle;
  \draw (0.55,0.96) -- (0.72,0.86);
  \draw (1.85,0.96) -- (1.68,0.86);
  \node[below, font=\small] at (1.2,-0.15) {sama kaki};
  \node[above, font=\small] at (1.2,2.1) {$A$};
  \node[below left, font=\small] at (0,0) {$B$};
  \node[below right, font=\small] at (2.4,0) {$C$};
  % sama sisi
  \begin{scope}[xshift=4.6cm]
  \draw[thick] (0,0) -- (2.4,0) -- (1.2,2.08) -- cycle;
  \draw (1.13,0.05) -- (1.27,-0.05);
  \draw (0.55,0.99) -- (0.72,0.89);
  \draw (1.85,0.99) -- (1.68,0.89);
  \node[below, font=\small] at (1.2,-0.25) {sama sisi};
  \end{scope}
  % siku-siku
  \begin{scope}[xshift=9.2cm]
  \draw[thick] (0,0) -- (2.6,0) -- (0,1.8) -- cycle;
  \draw[thin] (0.3,0) -- (0.3,0.3) -- (0,0.3);
  \node[below, font=\small] at (1.3,-0.15) {siku-siku};
  \node[omThm, font=\small, rotate=-34] at (1.5,1.05) {sisi miring};
  \end{scope}
\end{tikzpicture}

{\small Tiga segitiga istimewa. Garis kecil menandai sisi yang sama
panjang; persegi kecil menandai \omterm{def:g3:shapes:rightangle}{sudut siku-siku}.}
\end{omfigure}

\begin{method}[Melukis segitiga dari ketiga sisinya]\label{met:g6:shapes:construct}
Untuk melukis $ABC$ dengan $AB = 5$ cm, $AC = 4$ cm, $BC = 3$ cm:
\begin{enumerate}
  \item gambarlah \omterm{def:g6:lines:objects}{ruas garis} $[AB]$ sepanjang $5$ cm dengan penggaris;
  \item gambarlah busur lingkaran berpusat $A$ dengan \omterm{def:g3:shapes:shapes}{jari-jari} $4$ cm
        (semua titik yang berjarak $4$ cm dari $A$);
  \item gambarlah busur lingkaran berpusat $B$ dengan \omterm{def:g3:shapes:shapes}{jari-jari} $3$ cm;
  \item titik $C$ adalah tempat kedua busur itu berpotongan;
        hubungkan ia ke $A$ dan $B$.
\end{enumerate}
\end{method}

\begin{omfigure}
\begin{tikzpicture}[scale=1.0]
  \coordinate (A) at (0,0);
  \coordinate (B) at (5,0);
  \coordinate (C) at (3.2,2.4);
  \draw[thick] (A) -- (B);
  \draw[thick, omDef] (A) -- (C);
  \draw[thick, omThm] (B) -- (C);
  \draw[omDef, thin] (C) arc[start angle=36.9, delta angle=22, radius=4];
  \draw[omDef, thin] (C) arc[start angle=36.9, delta angle=-22, radius=4];
  \draw[omThm, thin] (C) arc[start angle=126.9, delta angle=22, radius=3];
  \draw[omThm, thin] (C) arc[start angle=126.9, delta angle=-22, radius=3];
  \fill (A) circle (1.5pt) node[below left, font=\small] {$A$};
  \fill (B) circle (1.5pt) node[below right, font=\small] {$B$};
  \fill (C) circle (1.5pt) node[above, font=\small] {$C$};
  \node[omDef, font=\small] at (1.2,1.5) {$4$ cm};
  \node[omThm, font=\small] at (4.6,1.4) {$3$ cm};
  \node[below, font=\small] at (2.5,0) {$5$ cm};
\end{tikzpicture}

{\small Melukis segitiga dengan jangka: \omterm{def:g5:solids:def}{titik sudut} $C$ adalah titik
potong dua busur. (Pada \cref{ch:g7:triangles} kamu akan melihat kapan
lukisan semacam itu mungkin dilakukan.)}
\end{omfigure}

\section{Segi empat}

\begin{definition}[Segi empat istimewa]\label{def:g6:shapes:quadrilaterals}
\omterm{def:g4:geometry:polygon}{Segi empat} punya empat \omterm{def:g5:solids:def}{titik sudut} dan empat sisi. Ia disebut:
\begin{itemize}
  \item \emph{persegi panjang}\index{persegi panjang} jika keempat
        sudutnya siku-siku;
  \item \emph{belah ketupat}\index{belah ketupat} jika keempat sisinya
        sama panjang;
  \item \emph{persegi}\index{persegi} jika ia keduanya sekaligus:
        empat \omterm{def:g3:shapes:rightangle}{sudut siku-siku} \emph{dan} empat sisi sama panjang.
\end{itemize}
\end{definition}

\begin{proposition}[Sifat diagonalnya]\label{prop:g6:shapes:diagonals}
Pada \omterm{def:g4:geometry:polygon}{segi empat} ini, kedua diagonalnya saling memotong di titik tengah
bersamanya, dan lebih dari itu:
\begin{itemize}
  \item pada persegi panjang, diagonalnya \emph{sama panjang};
  \item pada \omterm{def:g6:shapes:quadrilaterals}{belah ketupat}, diagonalnya \emph{saling \omterm{def:g6:lines:perp}{tegak lurus}};
  \item pada persegi, keduanya sekaligus.
\end{itemize}
\end{proposition}
\admitted

\begin{omfigure}
\begin{tikzpicture}[scale=0.8]
  % persegi panjang, tanda sudut siku-siku digambar ke arah dalam di tiap pojok
  \draw[thick] (0,0) rectangle (3,1.9);
  \draw[omThm] (0,0) -- (3,1.9);
  \draw[omThm] (0,1.9) -- (3,0);
  \foreach \x/\y/\dx/\dy in {0/0/1/1, 3/0/-1/1, 0/1.9/1/-1, 3/1.9/-1/-1}
    \draw[thin] (\x+0.22*\dx,\y) -- (\x+0.22*\dx,\y+0.22*\dy)
      -- (\x,\y+0.22*\dy);
  \node[below, font=\small] at (1.5,-0.2) {persegi panjang};
  % belah ketupat
  \begin{scope}[xshift=5.1cm, yshift=0.95cm]
  \draw[thick] (0,0) -- (1.6,1.0) -- (3.2,0) -- (1.6,-1.0) -- cycle;
  \draw[omThm] (0,0) -- (3.2,0);
  \draw[omThm] (1.6,1.0) -- (1.6,-1.0);
  \draw[thin] (1.6,0) ++(0.18,0) -- ++(0,0.18) -- ++(-0.18,0);
  \node[below, font=\small] at (1.6,-1.2) {belah ketupat};
  \end{scope}
  % persegi, beserta tanda sudut siku-sikunya juga
  \begin{scope}[xshift=9.8cm]
  \draw[thick] (0,0) rectangle (1.9,1.9);
  \draw[omThm] (0,0) -- (1.9,1.9);
  \draw[omThm] (0,1.9) -- (1.9,0);
  \foreach \x/\y/\dx/\dy in {0/0/1/1, 1.9/0/-1/1, 0/1.9/1/-1,
                             1.9/1.9/-1/-1}
    \draw[thin] (\x+0.22*\dx,\y) -- (\x+0.22*\dx,\y+0.22*\dy)
      -- (\x,\y+0.22*\dy);
  \node[below, font=\small] at (0.95,-0.2) {persegi};
  \end{scope}
\end{tikzpicture}

{\small Tiga \omterm{def:g4:geometry:polygon}{segi empat} istimewa dengan diagonalnya (warna merah).
Persegi mewarisi sifat persegi panjang sekaligus sifat \omterm{def:g6:shapes:quadrilaterals}{belah ketupat}.}
\end{omfigure}

\begin{remark}
Persegi itu otomatis merupakan persegi panjang (ia punya empat \omterm{def:g3:shapes:rightangle}{sudut
siku-siku}) dan otomatis merupakan \omterm{def:g6:shapes:quadrilaterals}{belah ketupat} (empat sisi sama
panjang). Jadi ``persegi panjang'' tidak berarti ``bukan persegi''!
Dalam matematika, bangun yang lebih istimewa termasuk ke dalam
keluarga yang lebih umum.
\end{remark}

\begin{example}[Mengenali bangun dari sifatnya]\label{ex:g6:shapes:recognize}
Sebuah \omterm{def:g4:geometry:polygon}{segi empat} punya empat sisi sepanjang $4$ cm dan satu \omterm{def:g3:shapes:rightangle}{sudut
siku-siku}. Bangun apa itu? Empat sisi sama panjang membuatnya \omterm{def:g6:shapes:quadrilaterals}{belah
ketupat}. Pada \omterm{def:g6:shapes:quadrilaterals}{belah ketupat}, sudut yang berhadapan sama besar dan
keempat sudutnya berjumlah $360^\circ$; satu \omterm{def:g3:shapes:rightangle}{sudut siku-siku} lalu
memaksa keempatnya menjadi siku-siku (hal ini akan dibuktikan bersama
jajargenjang pada \cref{ch:g7:central}). Jadi bangunnya persegi.
\end{example}

\begin{method}[Program melukis]\label{met:g6:shapes:program}
\emph{Program melukis} adalah daftar langkah yang dapat diikuti siapa
pun untuk menggambar ulang sebuah gambar dengan tepat. Tulislah satu
langkah per baris, namailah setiap titik ketika ia dibuat, dan
pakailah hanya: ``gambarlah \omterm{def:g6:lines:objects}{ruas garis} / garis / lingkaran\dots'',
``letakkan titiknya di tempat \dots\ berpotongan''. Ujilah programmu
dengan mengikutinya secara harfiah, kata demi kata.
\end{method}

\begin{example}\label{ex:g6:shapes:program}
Program untuk persegi $ABCD$ dengan sisi $4$ cm:
\begin{enumerate}
  \item Gambarlah \omterm{def:g6:lines:objects}{ruas garis} $[AB]$ dengan $AB = 4$ cm.
  \item Gambarlah \omterm{def:g6:lines:perp}{garis tegak lurus} terhadap $(AB)$ di $A$; padanya,
        letakkan $D$ pada jarak $4$ cm dari $A$ (pilihlah satu sisi
        garis lalu tetaplah di sana).
  \item Gambarlah \omterm{def:g6:lines:perp}{garis tegak lurus} terhadap $(AB)$ di $B$; padanya,
        letakkan $C$ pada jarak $4$ cm dari $B$, di sisi yang sama
        dengan $D$.
  \item Hubungkan $C$ dan $D$.
\end{enumerate}
Periksalah dengan penggaris: $DC$ semestinya juga berukuran $4$ cm.
\end{example}

\section{Latihan}

\begin{exercise}[$\star$]\label{exo:g6:shapes:1}
Lukislah segitiga $ABC$ dengan $AB = 6$ cm, $AC = 5$ cm dan
$BC = 4$ cm, mengikuti \cref{met:g6:shapes:construct}.
\end{exercise}

\begin{exercise}[$\star$]\label{exo:g6:shapes:2}
Lukislah \omterm{def:g6:shapes:triangles}{segitiga sama kaki} $DEF$ dengan $DE = DF = 5$ cm dan
$EF = 3$ cm. Dua sudut mana pada segitiga ini yang tampak sama besar?
Ukurlah untuk memeriksanya.
\end{exercise}

\begin{exercise}[$\star$]\label{exo:g6:shapes:3}
Lukislah \omterm{def:g6:shapes:triangles}{segitiga sama sisi} dengan sisi $4$ cm hanya memakai
penggaris dan jangka. Ukurlah salah satu sudutnya.
\end{exercise}

\begin{exercise}[$\star$]\label{exo:g6:shapes:4}
Golongkan setiap segitiga dari datanya (sama kaki, sama sisi,
siku-siku --- beberapa jawaban dapat berlaku sekaligus):
(a) sisi $7$, $7$, $7$; \
(b) sisi $5$, $5$, $8$; \
(c) sudut $90^\circ$, $45^\circ$, $45^\circ$; \
(d) sudut $60^\circ$, $60^\circ$, $60^\circ$.
\end{exercise}

\begin{exercise}[$\star$]\label{exo:g6:shapes:5}
Pada \omterm{def:g6:shapes:triangles}{segitiga siku-siku}, sisi mana yang menjadi sisi miring?
Gambarlah segitiga $MNP$ yang siku-siku di $N$, lalu sebutkan sisi
miringnya. \omterm{def:g5:solids:def}{Titik sudut} mana yang tidak pernah disentuh sisi miring itu?
\end{exercise}

\begin{exercise}[$\star$]\label{exo:g6:shapes:6}
Benar atau salah? Berilah alasan untuk tiap jawaban dalam satu kalimat.
\begin{enumerate}
  \item Setiap persegi adalah persegi panjang.
  \item Setiap persegi panjang adalah persegi.
  \item Setiap persegi adalah \omterm{def:g6:shapes:quadrilaterals}{belah ketupat}.
  \item \omterm{def:g6:shapes:quadrilaterals}{Belah ketupat} yang punya satu \omterm{def:g3:shapes:rightangle}{sudut siku-siku} adalah persegi.
\end{enumerate}
\end{exercise}

\begin{exercise}[$\star$]\label{exo:g6:shapes:7}
Lukislah persegi panjang $ABCD$ dengan $AB = 5$ cm dan $BC = 3$ cm.
Ukurlah kedua diagonalnya: apa yang kamu amati?
\end{exercise}

\begin{exercise}[$\star$]\label{exo:g6:shapes:8}
Lukislah \omterm{def:g6:shapes:quadrilaterals}{belah ketupat} yang diagonalnya $6$ cm dan $4$ cm. (Pakailah
\cref{prop:g6:shapes:diagonals}: gambarlah kedua diagonalnya dulu ---
saling \omterm{def:g6:lines:perp}{tegak lurus}, berpotongan di titik tengahnya --- lalu hubungkan
keempat ujungnya.)
\end{exercise}

\begin{exercise}[$\star\star$]\label{exo:g6:shapes:9}
Tulislah program melukis untuk \omterm{def:g6:shapes:triangles}{segitiga sama kaki} dengan alas $4$ cm
dan kaki $6$ cm, dengan gaya \cref{ex:g6:shapes:program}. Tukarkanlah
dengan teman sekelas: dapatkah ia mengikutinya tanpa melihat gambarmu?
\end{exercise}

\begin{exercise}[$\star\star$]\label{exo:g6:shapes:10}
Gambarlah lingkaran berpusat $O$ dan salah satu \omterm{def:g4:geometry:circle}{diameternya} $[AB]$.
Letakkan sembarang titik $M$ pada lingkaran lalu gambarlah segitiga
$AMB$. Ukurlah sudut $\widehat{AMB}$. Geserlah $M$ lalu ukur lagi. Apa
dugaanmu? (Fakta yang mengejutkan ini dibuktikan pada \cref{ch:g8:cosine}.)
\end{exercise}

\begin{exercise}[$\star\star$]\label{exo:g6:shapes:11}
Sebuah \omterm{def:g4:geometry:polygon}{segi empat} punya diagonal yang sama panjang dan saling
memotong di titik tengah bersamanya. Pilahlah bangun apa ia pastinya,
dengan memakai \cref{prop:g6:shapes:diagonals}: persegi panjang, \omterm{def:g6:shapes:quadrilaterals}{belah
ketupat} atau persegi? (Hati-hati: hanya satu sifat yang diberikan.)
\end{exercise}

\begin{exercise}[$\star\star\star$]\label{exo:g6:shapes:12}
Pada petak titik $4 \times 4$, berapa persegi berbeda yang dapat kamu
gambar dengan pojoknya pada titik petak --- termasuk juga persegi yang
``miring''? Gambarlah sedikitnya satu persegi miring lalu jelaskan
mengapa sisinya sama panjang.
\end{exercise}

\section{Soal: Apa yang diketahui oleh diagonalnya}

\begin{problem}[{Soal akhir pekan --- membangun kembali setiap segi
empat istimewa hanya dari diagonalnya, dan membilang diagonal segi
banyak mana pun}]
\label{pb:g6:shapes:1}
\Cref{prop:g6:shapes:diagonals} membaca sebuah bangun lalu
menceritakan diagonalnya. Soal ini memainkan permainan itu
\emph{secara terbalik}: diberikan hanya dua \omterm{def:g6:lines:objects}{ruas garis} yang
berpotongan --- calon diagonalnya --- bangun apa yang muncul kalau
keempat ujungnya dihubungkan? Kamu akan menyusun ``kamus diagonal''
yang lengkap, memahami \emph{mengapa} lukisan \omterm{def:g6:shapes:quadrilaterals}{belah ketupat} pada
\cref{exo:g6:shapes:8} berhasil, bertemu bangun yang belum punya
nama, lalu menutupnya dengan membilang diagonal \omterm{def:g4:geometry:polygon}{segi banyak} yang
sisinya jauh lebih dari empat.

\medskip
\noindent\textbf{Bagian I --- Kamus diagonal.}
Untuk setiap pertanyaan, gambarlah kedua \omterm{def:g6:lines:objects}{ruas garis} seperti yang
dijelaskan, hubungkan keempat ujungnya berurutan, lalu kenali \omterm{def:g4:geometry:polygon}{segi
empat} yang kamu peroleh.
\begin{enumerate}
  \item Dua \omterm{def:g6:lines:objects}{ruas garis} yang masing-masing $6$ cm, saling \omterm{def:g6:lines:perp}{tegak
        lurus}, berpotongan di titik tengah bersamanya. Ukurlah keempat sisi
        dan keempat sudut bangun yang dihasilkan: bangun apa itu?
  \item Dua \omterm{def:g6:lines:objects}{ruas garis} $8$ cm dan $5$ cm, saling \omterm{def:g6:lines:perp}{tegak lurus},
        berpotongan di titik tengah bersamanya. Bangun apa yang
        muncul (\cref{exo:g6:shapes:8})?
  \item Dua \omterm{def:g6:lines:objects}{ruas garis} yang masing-masing $7$ cm --- sama panjang tetapi
        \emph{tidak} saling \omterm{def:g6:lines:perp}{tegak lurus} --- berpotongan di titik
        tengah bersamanya. Ukurlah sisi dan sudutnya: bangun apa
        dugaanmu?
  \item Dua \omterm{def:g6:lines:objects}{ruas garis} $8$ cm dan $5$ cm, tidak saling \omterm{def:g6:lines:perp}{tegak lurus},
        berpotongan di titik tengah bersamanya. Bangun yang kamu
        peroleh bukan salah satu dari ketiga bangun pada
        \cref{def:g6:shapes:quadrilaterals} --- jelaskan apa yang kamu
        lihat (sisi mana yang tampak sama panjang? \omterm{def:g6:lines:perp}{sejajar}?). \omterm{def:g4:geometry:polygon}{Segi
        empat} ini menerima namanya pada \cref{ch:g7:central}.
  \item Rangkumlah pertanyaan 1--4 menjadi kamus dengan empat baris:
        diagonalnya \emph{sama panjang / \omterm{def:g6:lines:perp}{tegak lurus} / keduanya /
        bukan keduanya} (selalu berpotongan di titik tengahnya)
        $\rightarrow$ nama (atau gambaran) bangunnya.
\end{enumerate}

\medskip
\noindent\textbf{Bagian II --- Mengapa lukisannya berhasil.}
\begin{enumerate}[resume]
  \item Pada pertanyaan 2, kedua diagonalnya memotong bangun itu
        menjadi empat \omterm{def:g6:shapes:triangles}{segitiga siku-siku} di sekeliling titik
        potongnya. Berapa panjang kedua sisi siku-siku setiap
        segitiga itu? Jelaskan mengapa keempat segitiganya adalah
        salinan persis satu sama lain, lalu simpulkan mengapa keempat
        sisi bangunnya sama panjang --- rahasia \cref{exo:g6:shapes:8}.
  \item Simpulkan sebuah uji yang hanya memakai jangka: setelah
        membangun bangun apa pun dari diagonalnya, bagaimana kamu
        dapat memeriksa hanya dengan jangka (tanpa mengukur dengan
        penggaris) bahwa keempat sisinya sama panjang?
  \item Tulislah program melukis
        (\cref{met:g6:shapes:program}) untuk persegi yang
        \emph{diagonalnya} berukuran $6$ cm --- bukan sisinya.
  \item Pada persegi di pertanyaan 8, apakah sisinya lebih panjang
        atau lebih pendek daripada $3$ cm --- setengah diagonalnya?
        Ukurlah untuk mengetahuinya. Setengah jawabannya sudah dapat
        \emph{dibuktikan}: berjalan dari satu pojok ke pojok
        berikutnya lewat pusatnya adalah memutar sejauh $3 + 3$ cm,
        dan jalan lurus selalu lebih pendek daripada memutar --- jadi
        sisinya lebih pendek daripada $6$ cm. (Membuktikan bahwa ia
        juga mengalahkan $3$ cm, dan menghitungnya dengan tepat,
        adalah tugas \cref{ch:g8:pythagoras}.)
  \item Sebuah baris baru untuk kamusnya: gambarlah \omterm{def:g6:lines:objects}{ruas garis} $[AC]$
        sepanjang $8$ cm; padanya, tandailah titik $H$ berjarak $3$ cm
        dari $A$; melalui $H$ gambarlah ruas \omterm{def:g6:lines:perp}{garis tegak lurus} $[BD]$
        sepanjang $5$ cm dengan $H$ sebagai titik tengah \emph{ruas
        itu}. Hubungkan $ABCD$. Bangun ini disebut
        \emph{layang-layang}\index{layang-layang} (satu diagonalnya
        membagi dua diagonal yang lain, \omterm{def:g6:lines:perp}{tegak lurus}). Dengan memakai
        segitiga pojok seperti pada pertanyaan 6, carilah sisi mana
        pada \omterm{pb:g6:shapes:1}{layang-layang} yang sama panjang dengan sisi mana.
\end{enumerate}

\medskip
\noindent\textbf{Bagian III --- Membilang diagonal.}
\emph{Diagonal} sebuah \omterm{def:g4:geometry:polygon}{segi banyak} menghubungkan dua \omterm{def:g5:solids:def}{titik sudut} yang
bukan tetangga.
\begin{enumerate}[resume]
  \item Gambarlah \omterm{def:g4:geometry:polygon}{segi empat}, \omterm{def:g4:geometry:polygon}{segi lima} (lima \omterm{def:g5:solids:def}{titik sudut}) dan \omterm{def:g4:geometry:polygon}{segi
        enam} (enam \omterm{def:g5:solids:def}{titik sudut}), lalu gambarlah semua diagonalnya.
        Bilanglah: ada berapa pada masing-masing?
  \item Membilang tanpa menggambar, untuk \omterm{def:g4:geometry:polygon}{segi enam}: dari satu \omterm{def:g5:solids:def}{titik
        sudut}, ada berapa diagonal yang berangkat? (Ke berapa dari
        keenam \omterm{def:g5:solids:def}{titik sudut} itu ia \emph{tidak} dapat dihubungkan oleh
        diagonal?) Mengalikannya dengan keenam \omterm{def:g5:solids:def}{titik sudut} membilang
        setiap diagonal tepat dua kali --- mengapa? Perolehlah
        kembali hitunganmu pada pertanyaan 11. (Pembilangan ganda
        yang sama muncul pada jabat tangan di \cref{pb:g6:wholes:1}.)
  \item Pakailah cara pada pertanyaan 12 untuk membilang diagonal
        \omterm{def:g4:geometry:polygon}{segi banyak} dengan $10$ \omterm{def:g5:solids:def}{titik sudut}, lalu dengan $20$ \omterm{def:g5:solids:def}{titik
        sudut}.
  \item Sebuah \omterm{def:g4:geometry:polygon}{segi banyak} punya tepat $44$ diagonal. Berapa \omterm{def:g5:solids:def}{titik
        sudutnya}? (Cobalah cara pada pertanyaan 13 pada beberapa
        calon.)
  \item Ujilah caramu pada \omterm{def:g4:geometry:polygon}{segi banyak} yang paling kecil: berapa
        diagonal yang dipunyai segitiga, dan apa jawaban cara
        pembilangan pada pertanyaan 12? Tutuplah dengan satu kalimat
        tentang persamaan jabat tangan dan diagonal.

\end{enumerate}
\end{problem}
